% 【本文件写什么】impl 实现层：dummy
% - 定位：占位/链接桩实现，保证默认 feature 可编译。
% - crate 路径：os/components/wateros-vfs/vfs-impl/impl-dummy/src/
% - 如何实现对应 api-v0 trait：关键类型、状态机、数据结构
% - 算法或硬件交互流程，可用伪代码/流程图
% - 本 impl 启用的 Cargo [features] 与 arch feature，impl-riscv64 等
% - 与其它 impl 变体的差异、切换 feature 时的影响
% - 性能/内存假设与已知限制
% 【事实来源】
% - 上述 crate 源码与 Cargo.toml
% - os/feature-tree.txt 中指向本 impl 的 feature 链

\subsubsection{impl-dummy}

\code{DummyBackend}实现完整\code{VfsBackend}组合trait，但卷访问均返回\code{NotMounted}或\code{Unsupported}：\code{exists}/\code{metadata}/\code{read}在路径规范化后失败；\code{mount\_rw\_session}为\code{Unsupported}；\code{list\_dev\_nodes}为空；\code{VfsOpenOps}使用trait默认空实现。

路径规范化，即 \code{normalize\_absolute\_path}，仍可用，便于无\code{bridge-fs-api}配置下的api契约单测。当未启用\code{bridge-fs-api}时，\code{active\_impl::backend()}选中本后端。

不提供真实挂载、open或fd I/O；仅保证链接与\code{vfs::test()}中api/dummy烟囱通过。
